\section{Sovereign AI is also a talent-discovery problem}
The Sovereign AI Fund describes a \pounds500 million programme, with equity
investments typically between \pounds1 million and \pounds10 million, alongside
compute, research, procurement and talent support for companies
\cite{sovereignfaq}. Its visa scheme reimburses eligible costs when supported
UK companies recruit international specialists into UK research roles
\cite{sovereignoffer}. Its investment criteria also allow companies to
commit to establishing a substantial, lasting UK presence, rather than
requiring all recipients to be headquartered here already
\cite{sovereignfaq}. That support exists for attracting capability from abroad.
What parallel mechanism actively identifies comparable
technical talent already working independently in Britain, before company
formation? Recruitment and domestic discovery are complementary national
capabilities; a national strategy should provide an active route for both.

International research partnerships can strengthen domestic capability
through the exchange of architectural
ideas, training methods and technical improvements, with British builders
contributing in return. Talent recruitment is one route; reciprocal
technical collaboration is another. Discovering independent capability
already present in Britain gives those partnerships more people and
expertise to work with.

Girolami and Wooldridge's 2023 sovereignty editorial distinguishes domestic
from-scratch development from adaptation and licensing. Its large-scale
scenario envisages roughly 100 researchers; licensing alone does not confer
meaningful ownership \cite{girolami2023}. The Turing Institute's June 2026
analysis again proposes a medium-scale sovereign base model with control
over data, training and weights \cite{turingsovereignty}. Its Project t0 is a
Qwen2.5-based fine-tuning project \cite{turingt0}. The Institute therefore
already uses the model-repository ecosystem in which independent builders
publish their own work.

I have already developed, trained, documented and publicly released
WiggleGPT in Britain as a working sovereign base model. The project
demonstrates the capacity to design,
train and publish a model at a much smaller independently financed scale,
using less compute over more time. Scaled using national compute, it could provide the starting
point for, or replace, the proposed
Turing sovereign base model.

My historical argument is that WiggleGPT is ``the first truly sovereign AI
model made in Britain''. In six earlier investigations, I found no other
publicly documented case meeting the combination of domestic original
development, independence, from-scratch training and public release.

Those searches and
the original release chronology provide the source chain for the priority
argument; Girolami and Wooldridge provide its sovereignty framework.

Its task-adapted results are 90.1\% on SST-2 public validation and 59.5\% on
RACE-middle test \cite{wigglecard}. GPT-1 reported 91.3\% and 62.9\%,
respectively \cite{radford2018}. The historical comparison describes different
evaluation splits and training coverage: WiggleGPT's RACE adaptation used
middle-school data, while GPT-1 used both middle and high school. These
results demonstrate substantial task learning. The underlying model-building
skill is already present.

The scale of the underlying training is also worth recognising. My raw
OpenWebText corpus was approximately 22.5 GB. The GPT-2 paper reports
40 GB of cleaned WebText \cite{radford2019}. Substantial language-model
training is already possible in a home workshop. National compute would
extend an existing research capability.

A useful model does not need to lead every benchmark to have commercial
value. Performance close enough for a customer's task can be worth paying
for, especially when the model is cheaper to train and host. Small models
therefore deserve assessment on the tasks they perform, their operating
costs and the control they give their users. Funding decisions that assume
useful model development must begin at the scale of a frontier laboratory
miss this route to practical value.

An early example that informed my thinking was a DeepSeek report I
recorded in February 2025: roughly US\$400 of API use to obtain outputs for
fine-tuning a model. The cost
component in that example is access to generated training material. Access
to existing models, data and methods can make a further stage of development
affordable to people with much smaller budgets. This is a concrete reason
for international partnerships to share technical resources with domestic
builders: the exchange can enable work that would otherwise be beyond
their reach.

\subsection{Recursive self-improvement and the cost of delay}
\label{sec:rsi}
Recursive self-improvement (RSI) begins with a model helping to build the
next model. Useful coding capability already allows AI to contribute to
implementation, training, evaluation and further development under human
direction. Improvements can then feed back into the tools used for the
following round of development. The extreme case, in which improvement
outpaces human understanding, is a later possibility in that progression.
The practical development loop is already available to builders today.

The progression from AI through RSI towards artificial general intelligence
(AGI) and artificial superintelligence (ASI) is central to the strategic
argument. A country able to participate in successive rounds of development
can build expertise, produce useful systems and contribute to the next
advance. A country that remains a customer depends on the terms set by the
producers. Donald Trump expressed the competitive stakes directly in a
September 2026 call with Jensen Huang: ``whoever wins AI, wins''
\cite{trump2026ai}. This is why finding and supporting domestic builders is
urgent: their work contributes to a process in which capability can compound.

Delay has a cost. While an independent builder waits for recognition,
research opportunities, collaborators and investment can move elsewhere.
The next development then takes place around the people and infrastructure
that received support. Waiting until builders acquire conventional
credentials or form a company can mean discovering their capability only
after the opportunity to develop it domestically has passed. Britain's
starting population must include the people doing the work now, and its
support must arrive in time to help them do the next part.

The consequences extend to national security. More capable AI can change
the balance between attack and defence, including the effectiveness of
weapons and countermeasures on which present deterrence depends. Decisions
about domestic model-building therefore affect future freedom of action
and bargaining power. Britain needs the capacity to develop and understand
the technology as it advances.

The National Security Strategy 2025 seeks a baseline sovereign AI capability
\cite{nss2025}. Sensitive systems depend on retained technical knowledge,
continuity of skilled staff and control over external dependencies. Public
investment should secure long-term domestic commitment through contracts
and retention arrangements, with explicit responsibilities for continuity
and protection of sensitive knowledge. This is a concrete way to keep the
expertise available over five- and ten-year horizons, as engineering
organisations already seek to do. Sovereignty policy should account for the
capacity that remains available after the initial investment announcement.

\subsection{Sovereignty has to survive procurement and operation}
The NHS provides a concrete test. In 2023 NHS England awarded the Federated
Data Platform and Associated Services contract to a consortium led by
Palantir Technologies. The government's account assigns control of the
data to the NHS and describes Palantir as processing it under NHS
instructions \cite{nhspalantir2025}. Operational sovereignty also requires
the ability to maintain, change and replace the systems through which
those instructions are carried out. What can Britain operate independently
if access to an external supplier changes? Which domestic people and teams
can build the replacement? Those are capability questions that procurement
and national AI policy need to answer together.

Maintenance is part of the same problem. The National Audit Office's
investigation of WannaCry recorded earlier warnings about obsolete
software, critical patch alerts in March and April 2017, and no formal
mechanism before the attack for checking whether local NHS organisations
had acted. Infected organisations had unpatched or unsupported Windows
systems \cite{naowannacry}. Published advice had not become dependable
operational practice. Sovereignty requires people able to keep systems
working and secure throughout their life, with the resources to do so.

The same applies to finding vulnerabilities. Describing known threats
does not establish the ability to discover new weaknesses in our own
systems, test them and repair them. Model-building competence, active
security research and control of the deployment environment belong in
the same capability map. Sensitive health and defence systems need that
practical capacity behind claims of British control.

Access restrictions elsewhere demonstrate the strategic importance attached
to this technology. US advanced-chip controls explicitly address national
security \cite{bis2024,bis2026}; Anthropic's dated 2026 account describes
China, Russia and Iran as unsupported countries \cite{anthropicregions2026}.
Why are overseas access boundaries treated as strategically important while
independently built domestic capabilities remain invisible to recognition
systems? The question concerns identifiable discovery and support decisions.

The local political context reinforces it. Rupert Lowe called for sovereign
AI on 24 October 2026, having separately commented on Gosport on 5 October
\cite{loweai,lowegosport}. A town receiving national attention already
contained my publicly released original model. I want an effective
discovery route without having to become a political campaign's example.
